\scanpage{122}
% L1
$X$ Schéma compact, [i.e. quasi-compact et \inserted{\uline{topologiquement}} \uuline{séparé}]

% L2–6
(i) $X$ est tot. discontinu, \deleted{i.e. si $x,y\in X$, $x\ne y$, alors il existe une partie ouverte et fermée de $X$ qui contient $x$ et non $y$. En effet, soient $U,V$ des ouverts, $U\ni x$, $V\ni y$, $U\cap V=\varnothing$.} i.e. tt $x\in X$ a un syst. fond. de voisinages qui sont à la fois \uline{ouverts} et \uline{fermés}

% L7–8
[En effet, il y a un syst. fond. de voisinages ouverts quasi compacts]

% L9–11
(ii) $X$ \unsure{est|\deleted{est}} affine. En effet, \deleted{\unsure{on construit de proche en proche des ouverts affines non vides qui ne rencontrent pas}}\inserted{on recouvre $X$ par des ouverts affines en nb fini, soit $X_i$ $(1\leq i\leq n)$}. Ils sont ouverts et fermés. Soit
% L12–16
$X'_i=X_i-X_i\cap\bigcup_{j<i}X_j$, Alors $X'_i$ est une partie ouverte et fermée de $X_i$, donc affine puisque $X_i$ l'\unsure{est} (en vertu p. ex. du critère \unsure{coh.|ult.} \deleted{\unsure{,}} ou \unsure{n'imp. quoi|n'imp. quel}). \unsure{Comme|Donc} $X$ est \deleted{\unsure{est}} la somme des $X'_i$, il est affine.

% L17–19
(iii) Les parties constructibles de $X$ sont ses parties ouvertes et fermées. [\deleted{Donc}\inserted{\unsure{En dts|En ft|En effet}} $\operatorname{Cons} X\to X$ est un \deleted{\unsure{homéom.}]} \unsure{isom|iso.|bon}] Donc
% archive sheet number: 121
